\section{Related Work}\label{sec:related}
Early finite-key analyses of BB84 relied on Chernoff- and Hoeffding-type concentration bounds \cite{Hoeffding63,CaiScarani09}. Tomamichel \emph{et al.}\ \cite{Tomamichel12} obtained a tight, composable bound using the entropic uncertainty relation for smooth entropies \cite{TomamichelRenner11} together with Serfling's inequality for sampling without replacement \cite{Serfling74}. Hayashi and Tsurumaru \cite{HayashiTsurumaru12} gave a concise analysis based on different techniques, and Lim \emph{et al.}\ \cite{Lim14} extended finite-key bounds to decoy-state QKD. Imperfect sources, such as those exploitable by photon-number-splitting attacks, are treated in \cite{GLLP04,Brassard00} and are outside the scope of this paper. Our contribution is not a new security proof, but a self-contained quantitative study of how the established finite-key bound behaves for the statistics produced by IR attacks in noisy channels, together with an independent numerical validation.
